% Final
\chapter{Zadání uživatelského testování}\label{chapter:test-scenarios}

% Final
V této příloze lze nalézt zadání uživatelského testování použitelnosti aplikace.

% Final
\section{TS1 -- Registrace a zapsání kurzu}

% Final
\paragraph{Zadání} Otevřel jste aplikaci na učení se programování. Založte si v aplikaci účet, vyplňte úvodní dotazník a zapište si kurz \textit{HTML}. Při registraci použijte uživatelské jméno \textit{Test}, e-mail \textit{test@example.com} a heslo \textit{Test1234}.

% Final
\paragraph{Očekávaný průchod aplikací}

% Final
\begin{enumerate}
    \item Uživatel začíná na úvodní stránce aplikace.
    \item Uživatel se z úvodní stránky přesune na stránku registrace, vyplní registrační formulář a zvolí možnost zaregistrovat se. Aplikace přesune uživatele na stránku úvodního dotazníku.
    \item V dotazníku uživatel vyplní svoji úroveň programování, témata, která se chce učit, a cíl, kterého chce dosáhnout.
    \item Následně uživatel vybere jeden z doporučených kurzů.
    \item Nakonec uživatel zvolí svůj denní cíl učení.
    \item Aplikace oznámí uživateli, že je jeho účet připraven a uživatel zvolí možnost pokračovat, čímž se přesune na hlavní stránku aplikace.
\end{enumerate}

% Final
\section{TS2 -- Učení se dle studijního plánu}

% Final
\paragraph{Zadání} Otevřel jste aplikaci a chcete se začít učit vámi zapsaný kurz. Spusťte nadcházející lekci kurzu a lekci dokončete.

% Final
\paragraph{Očekávaný průchod aplikací}

% Final
\begin{enumerate}
    \item Uživatel začíná na hlavní stránce aplikace.
    \item Uživatel zvolí možnost spustit nadcházející lekci kurzu.
    \item Aplikace přesune uživatele na stránku učení, kde mu zobrazí látku nadcházející lekce.
    \item Uživatel splní všechna cvičení v lekci a aplikace mu oznámí, že byla lekce dokončena.
    \item Uživatel zvolí možnost pokračovat a aplikace ho přesměruje zpět na hlavní stránku.
\end{enumerate}

% Final
\section{TS3 -- Zapsání nového kurzu}

% Final
\paragraph{Zadání} Chcete si zapsat nový kurz \textit{CSS}. Zapište si tento kurz a spusťte v něm první lekci. Tuto lekci dokončete.

% Final
\paragraph{Očekávaný průchod aplikací}

% Final
\begin{enumerate}
    \item Uživatel začíná na hlavní stránce aplikace.
    \item Uživatel zvolí ikonu studijního plánu, čímž se přesune na stránku studijního plánu.
    \item Na stránce studijního plánu uživatel zvolí možnost přidat nový kurz, čímž se přesune na stránku s výběrem kurzů.
    \item Na stránce výběru kurzů uživatel zvolí kurz \textit{CSS}, čímž se přesune na stránku kurzu.
    \item Na stránce kurzu pak uživatel zvolí možnost začít kurz, čímž mu bude kurz zapsán a uživatel bude přesměrován na stránku učení první lekce kurzu.
    \item Uživatel splní všechna cvičení v lekci a aplikace mu oznámí, že byla lekce dokončena.
    \item Uživatel zvolí možnost pokračovat a aplikace ho přesměruje zpět na hlavní stránku.
\end{enumerate}

% Final
\section{TS4 -- Úprava studijního plánu}

% Final
\paragraph{Zadání} Máte zapsané kurzy \textit{HTML} a \textit{CSS}, které jste se nějakou dobu učili. Nyní se chcete zaměřit pouze na kurz CSS a z kurzu HTML se již nechcete učit nové lekce. Zařiďte, aby vám aplikace na hlavní stránce již nezobrazovala nové lekce z kurzu HTML, ale abyste nepřišli o svůj pokrok v kurzu HTML.

% Final
\paragraph{Očekávaný průchod aplikací}

% Final
\begin{enumerate}
    \item Uživatel začíná na hlavní stránce aplikace.
    \item Uživatel zvolí ikonu studijního plánu, čímž se přesune na stránku studijního plánu.
    \item Na stránce studijního plánu uživatel přesune kurz \textit{HTML} ze sekce Učení a opakování do sekce Pozastaveno. Aplikace mu tak na hlavní stránce přestane zobrazovat nové lekce z kurzu HTML.
\end{enumerate}

% Final
\section{TS5 -- Spuštění konkrétní lekce}

% Final
\paragraph{Zadání} V rámci kurzu HTML se chcete naučit látku z lekce \textit{Obrázky}. Spusťte tuto konkrétní lekci a dokončete ji.

% Final
\paragraph{Očekávaný průchod aplikací}

% Final
\begin{enumerate}
    \item Uživatel začíná na hlavní stránce aplikace.
    \item Uživatel zvolí ikonu studijního plánu, čímž se přesune na stránku studijního plánu.
    \item Na stránce studijního plánu uživatel zvolí kurz \textit{HTML}, čímž se přesune na stránku kurzu.
    \item Na stránce kurzu uživatel vybere první kapitolu kurzu, čímž se přesune na stránku dané kapitoly.
    \item Na stránce kapitoly kurzu uživatel najde lekci \textit{Obrázky} a zvolí možnost spuštění lekce.
    \item Uživatel splní všechna cvičení v lekci a aplikace mu oznámí, že byla lekce dokončena.
    \item Uživatel zvolí možnost pokračovat a aplikace ho přesměruje zpět na hlavní stránku.
\end{enumerate}

% Final
\section{TS6 -- Zakoupení položky a změna vzhledu avatara}

% Final
\paragraph{Zadání} Za učení jste získali virtuální měnu. Kupte si položku v obchodě a zajistěte, aby se tato položka zobrazovala na vašem profilovém obrázku.

% Final
\paragraph{Očekávaný průchod aplikací}

% Final
\begin{enumerate}
    \item Uživatel začíná na hlavní stránce aplikace.
    \item Uživatel zvolí ikonu profilu, čímž se přesune na stránku profilu.
    \item Na stránce profilu dále zvolí ikonu obchodu, čímž se přesune na stránku obchodu.
    \item Na stránce obchodu uživatel zvolí libovolnou položku a zvolí možnost koupit položku.
    \item Poté se uživatel kliknutím na tlačítko zpět přesune na stránku profilu.
    \item Na stránce profilu uživatel klikne na svého avatara, čímž se přesune na stránku přizpůsobení avatara.
    \item Na stránce přizpůsobení avatara zvolí zakoupenou položku, kterou chce přidat na svůj profilový obrázek a změnu uloží.
    \item Aplikace uživatele přesune na stránku profilu.
\end{enumerate}